\section{Solving the promised norm equations}
\label{sec:norms}

Division of divisor classes will require a function with a prescribed
norm.  This section constructs such a function when its existence is
known.  The construction uses linear algebra and extraction of roots in
finite reduced algebras, and its complexity is polynomial in the degree
of the cover.  In particular, the cover degree is not treated as a
constant.

\begin{theorem}[Promised norm solver]\label{norm:solve}
Let $k$ be an explicitly represented finite field of characteristic $p$,
let $q$ be an odd prime, and suppose that a primitive $q$th root of unity
$\zeta\in k$ and a generator of the full $q$-primary subgroup of $k^*$
are given.  Let $F\in k[X]$ be monic and square-free, of positive degree
$N$ divisible by $q$.  Put
\[
 K_0=k(X),\qquad
 \mathcal L=K_0[Y]/(Y^q-F),\qquad \rho(Y)=\zeta Y.
\]
Let $G\in K_0^*$, write $D$ for the maximum of its reduced numerator
and denominator degrees, and assume
\begin{equation}\label{norm:characteristic}
 p>(N+2D+1)^2.
\end{equation}
If $G$ belongs to $\operatorname{Norm}_{\mathcal L/K_0}(\mathcal L^*)$, a deterministic
algorithm constructs $h\in \mathcal L^*$ with
\[
 \prod_{j=0}^{q-1}\rho^j(h)=G.
\]
Its bit complexity, and the degrees of the rational coefficients of
$h$ in the basis $1,Y,\ldots,Y^{q-1}$, are bounded by fixed polynomials
in $q,N,D$, and $\log|k|$.  It uses no polynomial factorization oracle.
Its field-dependent branches use only zero tests, so the algorithm
admits the simultaneous execution of Lemma~\ref{ff:simultaneous}.
\end{theorem}

The characteristic bound is a convenient sufficient condition for the
scalar searches in Lemma~\ref{ff:unitroot}; the main algorithm satisfies
it with considerable room.  The roots of $F$ need not be supplied, and
the norm solver does not require $F$ to split over $k$.
The function automorphism $\rho$ is the inverse of the action $\sigma$
used in the divisor sections. Their powers run over the same cyclic
group, so they define the same norm.

We first turn the norm equation into the problem of constructing a
rank-one idempotent in an algebra.  We then obtain that idempotent from
spaces of matrices regular at the places of the rational function
field.  This use of intersections of maximal orders is closely related
to the explicit matrix-algebra algorithms of
Ivanyos--Kutas--R\'onyai~\cite{IKR18}.  Their general algorithm permits a
finite-field factorization oracle.  Here the cyclic presentation
allows the local orders to be written explicitly, and a fiber
calculation supplies the idempotent by linear algebra.  We give the
construction and its bounds in full.

\subsection{The cyclic algebra and its local matrix models}

The polynomial $Y^q-F$ is irreducible over $K_0$: a place at which $F$
has valuation one is Eisenstein for it.  Thus $\mathcal L$ is a field, cyclic of
degree $q$ over $K_0$.  Form the $K_0$-algebra
\begin{equation}\label{norm:cyclic-algebra}
 \mathcal E=\bigoplus_{b=0}^{q-1} \mathcal L V^b,
 \qquad V a=\rho(a)V\ (a\in \mathcal L),\qquad V^q=G.
\end{equation}
Its multiplication table in the basis
\[
 W_{ab}=Y^aV^b,\qquad 0\le a,b<q,
\]
is explicit.  The norm promise implies that $\mathcal E$ is a full
matrix algebra over $K_0$, but does not initially give us an
isomorphism.  Indeed, if $h_0$ is any norm solution, let $\mathcal L$ act on
itself by multiplication and let $V$ act by $h_0\rho$.  These operators
satisfy \eqref{norm:cyclic-algebra}.  Their linear independence follows
by applying a putative relation between the operators $\rho^b$ to
$1,Y,\ldots,Y^{q-1}$ and using the invertible matrix
$(\zeta^{bj})_{b,j}$.  Dimension then gives an isomorphism
\begin{equation}\label{norm:abstract-splitting}
 \mathcal E\simeq M_q(K_0).
\end{equation}
This argument establishes existence only; the algorithm never uses
$h_0$ or this unspecified isomorphism.

A \emph{rank-one idempotent} in $\mathcal E$ is an element $e$ with
$e^2=e$ whose image under \eqref{norm:abstract-splitting} has rank one.
The following observation explains why constructing such an element
solves the original norm equation.

\begin{lemma}\label{norm:idempotent-to-norm}
Given a rank-one idempotent $e\in\mathcal E$, linear algebra over
$K_0$ constructs a norm solution $h\in \mathcal L$.
\end{lemma}

\begin{proof}
The left ideal $\mathcal E e$ has dimension $q$ over $K_0$.  The map
$\mathcal L\to\mathcal E e$, $a\mapsto ae$, is injective, since any nonzero
$a\in \mathcal L$ is invertible.  Both sides have dimension $q$, so it is an
isomorphism.  In particular, there is a unique $h\in \mathcal L$ satisfying
\begin{equation}\label{norm:extract-h}
 Ve=he.
\end{equation}
Solve this equation for the $q$ coefficients of $h$ in the basis
$1,Y,\ldots,Y^{q-1}$, using the explicit $q^2$-element basis of
$\mathcal E$ to express both sides.  Repeatedly multiply the equation
on the left by $V$, using $Va=\rho(a)V$.  After $q$ steps this gives
\[
 G e=V^q e=\left(\prod_{j=0}^{q-1}\rho^j(h)\right)e.
\]
Since $e\ne0$ and the two coefficients are scalars in $K_0$, their
equality proves the norm equation.  In particular $h\ne0$.
\end{proof}

Our remaining task is to construct $e$.  We will impose local
integrality conditions on elements of $\mathcal E$, then find a
rank-one matrix in their values at infinity.

Let $v$ be a closed point of $\mathbb P^1_k$.  Write $K_v$ for the
completed rational function field, $R_v$ for its valuation ring,
$\kappa_v$ for its residue field, and $T$ for a uniformizer.  A
\emph{lattice} in a finite-dimensional $K_v$-space is a free $R_v$-module
of full rank.  An \emph{order} in $\mathcal E\otimes K_v$ is a lattice
that contains $1$ and is closed under multiplication; it is
\emph{maximal} if no strictly larger order contains it.  We construct
one maximal order $\mathcal A_v$ at every $v$.  Only the points dividing
$F$, the numerator or denominator of $G$, and the point at infinity
require nontrivial conditions.

We will repeatedly use the matrix trace $\operatorname{tr}$ on the
split algebra.  It does not depend on the choice of isomorphism in
\eqref{norm:abstract-splitting}.  Conjugation by $Y$ and $V$ shows that
every nonidentity basis monomial has trace zero, while
$\operatorname{tr}(1)=q$.  Consequently
\begin{equation}\label{norm:trace-coordinates}
 a=\sum_{a',b'} c_{a'b'}W_{a'b'}
 \quad\Longrightarrow\quad
 c_{ab}=q^{-1}\operatorname{tr}(aW_{ab}^{-1}).
\end{equation}
Here $q$ is invertible by \eqref{norm:characteristic}.

\paragraph{Unramified places with $q\mid v(G)$.}
Suppose $v$ is finite, $v(F)=0$, and $w=v(G)$ is divisible by $q$.
Set $\widetilde V=T^{-w/q}V$, and take the $R_v$-span of
$Y^a\widetilde V^b$.  Both $Y$ and $\widetilde V$ have unit $q$th
powers, so this span is an order and contains the inverses of all its
basis monomials.  It is maximal.  To see this, any larger order has
integral matrix traces: its regular left-action trace is integral and
equals $q$ times matrix trace.  Pairing an element of the larger order
with the inverse basis monomials in
\eqref{norm:trace-coordinates} therefore makes all its coordinates
integral.  It already belongs to the displayed span.  This description
applies at all finite places outside the indicated finite set, with
$w=0$.

\paragraph{Places ramified in $\mathcal L$.}
Suppose $v(F)=1$, and again put $w=v(G)$.  Since $q$ is odd,
\begin{equation}\label{norm:ramified-unit}
 V'=Y^{-w}V,\qquad (V')^q=u:=G/F^w\in R_v^*.
\end{equation}
The commutation factor in taking the $q$th power is a power of
$\zeta^{q(q-1)/2}=1$.  The residue of $u$ is a $q$th power in
$\kappa_v^*$.  Indeed the extension at $v$ is totally ramified of
degree $q$.  For any promised solution $h_0$ one has
$v_{\mathcal L}(h_0)=v(G)=w$, with $v_{\mathcal L}(Y)=1$.  Hence $h_0/Y^w$ is a unit, its
norm is $u$, and the residue of that norm is the $q$th power of its
residue.

Take a residue root and lift it by Hensel's method to $R\in R_v^*$
with $R^q=u$.  On column vectors $e_0,\ldots,e_{q-1}$ define
\begin{equation}\label{norm:ramified-matrices}
 V'e_j=\zeta^jR e_j,\qquad
 Ye_j=\begin{cases}e_{j+1},&j<q-1,\\F e_0,&j=q-1.\end{cases}
\end{equation}
These matrices satisfy $V'Y=\zeta YV'$, $Y^q=F$, and $(V')^q=u$.
The distinct diagonal entries of $V'$ give all diagonal matrix units
by interpolation, and the invertible shift $Y$ gives the remaining
matrix units over $K_v$.  Thus the model is an isomorphism to
$M_q(K_v)$.  Let $\mathcal A_v$ be the inverse image of $M_q(R_v)$.

\paragraph{Unramified places with $q\nmid v(G)$.}
Suppose $v(F)=0$ and $q\nmid w=v(G)$.  The residue of $F$ must be a
$q$th power.  Otherwise, since $\kappa_v$ contains $\mu_q$ and $q$ is
prime, $Y^q-\overline F$ is irreducible and defines an unramified
extension of degree $q$.  All its norm valuations are then divisible
by $q$, contradicting the promise on $G$.  Lift a residue root to
$S\in R_v^*$ with $S^q=F$.  Use the matrices
\begin{equation}\label{norm:unramified-matrices}
 Ye_j=\zeta^{-j}S e_j,\qquad
 Ve_j=\begin{cases}e_{j+1},&j<q-1,\\G e_0,&j=q-1.\end{cases}
\end{equation}
Now $VY=\zeta YV$, including at the wrap from $e_{q-1}$ to $e_0$.
The same matrix-unit argument gives an isomorphism to $M_q(K_v)$;
again take the inverse image of $M_q(R_v)$.

\paragraph{Infinity.}
Put $z=X^{-1}$ and $\widetilde Y=z^{N/q}Y$.  Then
\[
 \widetilde Y^q=z^NF(z^{-1})=:F_\infty(z),\qquad F_\infty(0)=1.
\]
There is a unique root $S\in k[[z]]$ of $S^q=F_\infty$ with $S(0)=1$.
Use \eqref{norm:unramified-matrices} with $\widetilde Y$ in place of
$Y$, and with the same weighted shift for $V$, whatever the valuation
of $G$ may be.  The inverse image of $M_q(k[[z]])$ is $\mathcal A_\infty$.
This model also specifies reduction to $M_q(k)$ at infinity.

The matrix orders just used are maximal.  For example, pairing with
matrix units shows that any element in a larger order has integral
matrix entries, using the integrality of matrix trace exactly as
above.  We have therefore specified maximal orders at every place,
using only residue roots whose existence follows from the norm
promise.

\subsection{Why one fiber contains a rank-one matrix}

Define the finite-dimensional $k$-algebra of global sections and its
evaluation image by
\begin{equation}\label{norm:sections-fiber}
 \mathcal S=\{a\in\mathcal E:a\in\mathcal A_v\text{ for every }v\},
 \qquad
 \mathcal B=\operatorname{ev}_\infty(\mathcal S)\subseteq M_q(k).
\end{equation}
We will compute both spaces by linear equations.  First we establish
the feature of $\mathcal B$ that will make this useful: up to change of
basis it is the algebra of all block upper triangular matrices.  Its
first block therefore contains the image of a rank-one idempotent.

Here is the geometric explanation, including the descent from
completed local orders.  Intersect $\mathcal A_v$ with the rational
algebra $\mathcal E$.  This is a full lattice over the ordinary local
discrete valuation ring at $v$, and its completion is
$\mathcal A_v$.  Indeed the completed lattice lies between two powers
of a rational uniformizer times a fixed rational lattice.  Rational
approximation to a basis, to sufficiently high precision, produces a
basis of the same completed lattice.  Equivalently, the finitely many
congruences between these two bounds descend to the ordinary local
ring.  Maximality descends as well, since a larger ordinary order
would yield a larger completed order.

For the following existence argument only, fix an isomorphism
\eqref{norm:abstract-splitting}.  An ordinary local maximal order in
$M_q(K_0)$ is the endomorphism ring of a lattice.  In fact, apply the
order to the standard local lattice $R^q$ and take its span $M$.
The order is a finite module and contains the identity, so $M$ is a
full lattice stable under it.  The containing order
$\operatorname{End}_R(M)$ must equal it by maximality.  At all but
finitely many places the standard rational basis matrices and their
inverse change of coordinates are integral; there we may use $M=R^q$.
The resulting lattices glue to a vector bundle $\mathcal V$ of rank
$q$ on $\mathbb P^1_k$.  Concretely, its local sections are rational
column vectors satisfying the specified lattice conditions.  A basis
of one local lattice gives a trivialization on a neighborhood after
excluding the finitely many poles and other exceptional points.
We have proved
\begin{equation}\label{norm:endomorphism-bundle}
 \mathcal S=H^0(\mathbb P^1_k,\operatorname{End}\mathcal V).
\end{equation}
The chosen matrix model at infinity identifies the endomorphism
fiber with $M_q(k)$.  Any two such identifications differ by
conjugation: an automorphism of a full matrix algebra sends matrix
units to matrix units, from which a change of column basis is
obtained.  Thus its particular choice will not affect the asserted
shape of $\mathcal B$.

The following form of the projective-line splitting theorem holds
over the original field, rather than only after extending its
constants; see Hazewinkel--Martin~\cite{HazewinkelMartin82}.

\begin{lemma}[Splitting on the projective line]\label{norm:bundle-split}
Every vector bundle on $\mathbb P^1_k$ is a direct sum of line bundles
$\mathcal O(d)$, with integer degrees $d$.
\end{lemma}

\begin{proof}
We recall a proof over the given field $k$.  Line bundles on
$\mathbb P^1_k$ are classified by their degree.  On the two standard
affine charts, the Laurent-polynomial computation of their
cohomology gives
\[
 H^0(\mathcal O(t))=0\ (t<0),\qquad
 H^1(\mathcal O(t))=0\ (t\ge-1).
\]
The first assertion is immediate from polynomial degrees; for the
second, the \v Cech quotient is spanned by the Laurent monomials
between the two chart ranges, and that range is empty for $t\ge-1$.

Choose the largest integer $d$ for which there is a nonzero map
$\mathcal O(d)\to\mathcal V$.  Such integers exist and are bounded
above: enclosing the finitely many defining lattices of $\mathcal V$
between twists of a fixed trivial bundle proves both claims.  A
nonzero map is generically injective.  Saturating its image would
increase its degree by the effective divisor of its zeros, so
maximality of $d$ implies that it is a line subbundle.  The quotient
$\mathcal Q$ is a vector bundle, and induction on rank writes it as
$\bigoplus_i\mathcal O(t_i)$.  Twist
\[
 0\longrightarrow\mathcal O(d)\longrightarrow\mathcal V
   \longrightarrow\mathcal Q\longrightarrow0
\]
by $-d-1$.  Maximality of $d$ gives
$H^0(\mathcal V(-d-1))=0$, and $H^1(\mathcal O(-1))=0$.
The cohomology sequence therefore gives
$H^0(\mathcal Q(-d-1))=0$, whence $t_i\le d$ for every $i$.
Finally, each extension group
\[
 \operatorname{Ext}^1(\mathcal O(t_i),\mathcal O(d))
   =H^1(\mathcal O(d-t_i))
\]
vanishes.  The sequence splits, completing the induction.
\end{proof}

Apply the lemma and order the summands by decreasing degree:
$\mathcal V=\bigoplus_{i=1}^q\mathcal O(d_i)$ with
$d_1\ge\cdots\ge d_q$.  The $(i,j)$ entry of a global endomorphism
is a section of $\mathcal O(d_i-d_j)$.  If $d_i<d_j$ it is zero;
otherwise its value at infinity is arbitrary, because a line bundle
of nonnegative degree has a section with any prescribed value at that
point.  Grouping the equal degrees proves the promised block upper
triangular description of $\mathcal B$.  This is an existence proof of
the shape, not an algorithm for computing the splitting of
$\mathcal V$.

\begin{lemma}[Extracting a rank-one fiber idempotent]
\label{norm:fiber-idempotent}
Given a basis of a subalgebra $\mathcal B\subseteq M_q(k)$ that is
conjugate to a full block upper triangular algebra, linear algebra
over $k$ constructs a rank-one idempotent in $\mathcal B$.
\end{lemma}

\begin{proof}
Compute the kernel $\mathcal N$ of the bilinear form
\[
 (a,b)\longmapsto\operatorname{tr}(ab),\qquad a,b\in\mathcal B.
\]
In block upper triangular coordinates, a product has diagonal blocks
equal to the products of the corresponding diagonal blocks.  The
ordinary trace pairing on a full matrix algebra is nondegenerate in
every characteristic, as is seen by pairing its matrix units.
Hence $\mathcal N$ is exactly the strictly block upper triangular
part.  The common kernel
\[
 U=\bigcap_{n\in\mathcal N}\ker n
\]
is the column space of the first block; if there is only one block,
$\mathcal N=0$ and $U=k^q$.  Indeed every later block has arbitrary
maps to the first block, which exclude any vector with a nonzero
later component from this common kernel.

Choose $0\ne v\in U$ and a row $w$ with $wv=1$.  Then $a_0=vw$
has rank one and satisfies $a_0^2=a_0$.  In block coordinates it has
nonzero entries only in the first block row, so it belongs to
$\mathcal B$.  The computation requires only kernels and one scalar
normalization; it does not require finding the conjugating matrix.
\end{proof}

\begin{figure}[tb]
\centering
\[
 \mathcal B\sim
 \begin{pmatrix} *&*&*\\0&*&*\\0&0&*\end{pmatrix}
 \quad\xrightarrow{\ \ker\operatorname{tr}(ab)\ }\quad
 \mathcal N\sim
 \begin{pmatrix}0&*&*\\0&0&*\\0&0&0\end{pmatrix},
 \qquad
 U\sim\begin{pmatrix}k^{r_1}\\0\\0\end{pmatrix}.
\]
\caption{The fiber calculation with three distinct degree blocks.
Each star denotes an arbitrary matrix block, and the diagonal block
sizes are $r_1,r_2,r_3$.  The common kernel $U$ of the trace-pairing
kernel is the first block column space.  Every matrix $vw$ with
$v\in U$ belongs to $\mathcal B$; choosing $wv=1$ gives the rank-one
idempotent.  The calculation works in the original coordinates.}
\label{fig:norm-fiber}
\end{figure}

Figure~\ref{fig:norm-fiber} isolates the part of the vector-bundle
argument that is actually computed.  The splitting theorem guarantees
the triangular shape; the trace pairing finds enough of that shape
to construct a rank-one matrix without a factorization or a change
of basis.

\subsection{From the fiber idempotent to a norm solution}

Lift the matrix $a_0$ of Lemma~\ref{norm:fiber-idempotent} to any
$b\in\mathcal S$ with $\operatorname{ev}_\infty(b)=a_0$.  This is a
linear solve once \eqref{norm:sections-fiber} has been computed.  The
coefficients of the generic characteristic polynomial of $b$ are
regular at every point, since $b$ is a global endomorphism of
$\mathcal V$.  A rational function regular on $\mathbb P^1_k$ is
constant.  Evaluating at infinity consequently gives
\[
 \det(T-b)=T^{q-1}(T-1).
\]
The primary decomposition for the relatively prime factors
$T^{q-1}$ and $T-1$ shows that
\begin{equation}\label{norm:generic-idempotent}
 e=b^q
\end{equation}
is the projection onto the one-dimensional $1$-eigenspace: it is zero
on the nilpotent part and the identity on that eigenspace.  Thus $e$
is a rank-one idempotent in $\mathcal E$.

Lemma~\ref{norm:idempotent-to-norm} now recovers a norm solution by
solving $Ve=he$ over $K_0$.

It remains to make the spaces in \eqref{norm:sections-fiber}
computable with polynomially many coefficients.  This is where the
explicit local models, rather than an unspecified splitting of the
cyclic algebra, are essential.

\subsection{Computing all sections by finite linear systems}
\label{norm:section-computation}

Write $G=A/B$ with coprime nonzero polynomials $A,B\in k[X]$.
Square-free decomposition and gcds partition the finite places
dividing $FAB$ into pairwise coprime monic square-free polynomials
$H_i$ such that the two values
\[
 \bigl(v(F),v(G)\bigr)=(\epsilon_i,w_i)
\]
are constant over the irreducible factors of each $H_i$.
No irreducible factorization is involved: decompose the multiplicity
layers of $A$ and $B$ and refine their intersections with each other
and with $F$ by gcds.  All these polynomials have degree less than
$p$ by \eqref{norm:characteristic}, so derivative-based square-free
decomposition has no inseparable residual case.  Put
\begin{equation}\label{norm:H-M}
 H=\prod_i H_i,\qquad d_H=\deg H\le N+2D,
 \qquad M=10q^2(N+D+1).
\end{equation}

At every place in group $i$ we may use the same rational uniformizer
$T=H_i$.  Arithmetic modulo $H_i^s$ carries out, simultaneously, all
the required arithmetic modulo the $s$th power of these local
uniformizers.  The Chinese remainder theorem explains this statement
without requiring the factors of $H_i$ to be computed.  When a local
model requires a residue root, its residue is a unit modulo $H_i$
and is a $q$th power in every residue-field component by the norm
promise.  Lemma~\ref{ff:unitroot} constructs the root in $k[X]/H_i$.
The characteristic bound supplies more than $d_H^2$ distinct prime
field scalars for that lemma.  Hensel lifting then works modulo powers
of $H_i$, since the derivative $qR^{q-1}$ or $qS^{q-1}$ is a unit.

\begin{lemma}[Bounds for the local models]\label{norm:local-bounds}
At a place with one of the matrix models above, every nonzero entry
of the matrices of $W_{ab}$ and $W_{ab}^{-1}$ has valuation between
$-M$ and $M$.  At a place with the scaled-basis order, the exponents
changing between that basis and $W_{ab}$ have absolute value at most
$M$.  Consequently, the standard coefficients of every local-order
element have valuation at least $-M$ at the exceptional places and
at infinity.  At all other finite places those coefficients are
integral.
\end{lemma}

\begin{proof}
At a ramified finite place, let $D_R$ denote the unit diagonal matrix
of $V'$.  Then $V=Y^wD_R$, and commutation gives
\[
 W_{ab}=\zeta^{wb(b-1)/2}Y^{a+wb}D_R^b.
\]
For any integer $t$, the nonzero entry of $Y^t$ in column $j$ has
valuation $\lfloor(j+t)/q\rfloor$.  Since $0\le a,b<q$ and
$|w|\le D$, these valuations, and those of the inverse matrices,
have absolute value at most $D+2$.  In the unramified matrix model,
$Y$ is a unit diagonal matrix and each entry of $V^b$, for $b<q$,
has valuation either $0$ or $w$; inverse entries have the opposite
valuations.  At infinity the extra factor $z^{-aN/q}$ in $Y^a$
gives the bound $N+D$.  These bounds are smaller than $M$.
In the scaled-basis case, the change is multiplication by
$T^{\pm bw/q}$, whose exponent has absolute value at most $D$.

For an element in a matrix order, formula
\eqref{norm:trace-coordinates} and the bound on the inverse basis
matrices give the coefficient bound $-M$.  The scaled-basis assertion
gives the same conclusion in the remaining case.  At an ordinary
finite place the order is the unscaled standard span, so the
coefficients are integral there.
\end{proof}

It follows that every global section is included in the ansatz
\begin{equation}\label{norm:ansatz}
 b=\sum_{a,b'=0}^{q-1}c_{ab'}W_{ab'},\qquad
 c_{ab'}=\frac{P_{ab'}(X)}{H(X)^M},\qquad
 \deg P_{ab'}\le M(d_H+1).
\end{equation}
Indeed the denominator clears every finite pole, and the bound at
infinity says precisely that its numerator degree exceeds its
denominator degree by at most $M$.  The polynomial coefficients in
\eqref{norm:ansatz} are our scalar unknowns.

For each finite group impose membership in $\mathcal A_v$ as follows.
In a matrix model compute, modulo $T^{2M+1}$, the integral quantities
$T^M c_{ab'}$ and every matrix entry of $T^M W_{ab'}$.
Their products and sums give $T^{2M}$ times the matrix of the section.
The section is integral exactly when all coefficients below order
$2M$ vanish.  These are linear conditions over $k$.  In a
scaled-basis order use the coordinates in that basis in place of
matrix entries and impose the identical divisibility condition.
At infinity make the same computation with $T=z$ and the prescribed
matrix model.  The coefficient of $z^{2M}$, after the lower ones
vanish, is exactly the evaluated matrix in $M_q(k)$.

For completeness, the truncated entries are computed without an
infinite-series oracle.  Each nonzero basis entry has the form
$T^e u$ for a known exponent $e$ with $|e|\le M$ and a unit $u$ given
by the diagonal and weighted-shift formulas.  Lift the required unit
roots to precision $3M+3$ by the simple-root Hensel iteration.  After
multiplication by $T^M$, this precision more than supplies all terms
through $T^{2M}$.  For the rational coefficients of
\eqref{norm:ansatz}, cancelling $T^M$ from the denominator leaves a
unit denominator at the group, inverted by the polynomial Euclidean
algorithm.  At infinity replace $X$ by $z^{-1}$ and cancel the known
powers of $z$.  The same finite-precision computation then applies.

Solving these linear equations gives a basis of $\mathcal S$ together
with the evaluated matrix of every basis vector.  Taking their span
gives $\mathcal B$ and retains a linear map with which to lift its
elements.  The ansatz contains every section by
Lemma~\ref{norm:local-bounds}; the imposed conditions are necessary
and sufficient at every exceptional place, and elsewhere the ansatz
is already integral.  Thus the computed spaces are exactly those in
\eqref{norm:sections-fiber}.

\subsection{Algorithm and polynomial bounds}

The complete procedure can now be stated without any existential
computational step.

\begin{enumerate}
\item Form the multiplication table of $\mathcal E$ in the basis
      $Y^aV^b$.  Compute the square-free groups $H_i$ and the local
      models, obtaining their residue roots by
      Lemma~\ref{ff:unitroot} and their needed precision by Hensel
      lifting.
\item Use \eqref{norm:ansatz} and the divisibility conditions above
      to compute $\mathcal S$ and its evaluation map to $M_q(k)$.
\item In the image $\mathcal B$, compute the trace-pairing kernel,
      its common column kernel, and a rank-one idempotent $a_0=vw$
      as in Lemma~\ref{norm:fiber-idempotent}.
\item Lift $a_0$ to $b\in\mathcal S$ by a linear solve and compute
      $e=b^q$ with the multiplication table.
\item Solve $Ve=he$ for $h\in \mathcal L$ over $k(X)$ and return $h$.
\end{enumerate}

Every step is deterministic after fixing a basis order and the first
available pivot in each linear solve.  The preceding arguments prove
existence of the required roots, a nonzero vector in the common
kernel, a lift of $a_0$, and a solution of the last system; therefore
the procedure terminates on every promised input.

We record bounds to ensure that this is also a uniform algorithm.
There are at most $d_H$ groups, and
\[
 d_H\le N+2D,\qquad M=10q^2(N+D+1).
\]
The number of unknowns in \eqref{norm:ansatz} is
$q^2(M(d_H+1)+1)$.  The total number of scalar local conditions is at
most a constant times $q^2M(d_H+1)$, since precision at a group of
degree $d_i$ contributes $O(Md_i)$ coefficients.  All quotient
polynomial rings have dimension at most $(3M+3)d_H$ over $k$.
These are fixed polynomial bounds in $q,N,D$.  The finite-field root
calculations also have polynomial complexity by
Lemma~\ref{ff:unitroot}; the bit lengths of their powering exponents
are $O(d_H\log|k|)$.

For explicit denominator control, compute $b^q$ by successive
multiplications by $b$.  Write $P_*=M(d_H+1)$.  The coefficients of
$b^s$ have the common denominator
\[
 H^{Ms}B^{s-1}
\]
and numerator degrees at most
$sP_*+(s-1)(N+D)$.  Indeed a product of two standard basis monomials
introduces at most one factor of $F$ and one factor of $G=A/B$.
Passing to the displayed common denominator at the next multiplication
therefore adds at most $P_*+N+D$ to numerator degree, including every
addition of terms.  This proves the assertion by induction for
$1\le s\le q$.  The final system $Ve=he$ has the same common
denominator with at most one extra factor of $B$; its cleared entries
have polynomial degree.  A solution of this consistent rank-$q$
system is expressed by ratios of minors, of degrees at most $q$
times the cleared entry degree.  Fraction-reducing elimination or
these minor formulas therefore bounds the output degree and the
number of field operations by fixed polynomials.  Standard dense
polynomial arithmetic and Gaussian elimination convert these bounds
to a fixed polynomial in $q,N,D,\log|k|$ bit operations.  The more
generous common bound needed by the full factorization algorithm is
recorded in Proposition~\ref{complexity:bound}.

Finally, polynomial gcds, Hensel lifting, and linear elimination use
only field arithmetic and zero tests.  Integer choices such as
degrees and multiplicities are determined by those tests; the
valuation groups are formed by polynomial gcds rather than by an
unknown factorization.  The finite-field root routine has the same
property.  Hence all steps are compatible with simultaneous
execution, proving the last assertion of Theorem~\ref{norm:solve}
and completing its proof.
